\section{Domain libraries built on the common kernel}
\label{sec:domain-libraries}

The framework is useful only if domain proofs share substantial infrastructure.
This section gives the intended library boundaries and the preferred statement
forms for the main areas represented in the repository.

\subsection{Exact finite algebra}

Polynomial and matrix identities should be proved first over coded coefficient
rings by finite sums and products.  A polynomial is a bounded coefficient list;
addition, multiplication, derivative, composition, and evaluation are
primitive-recursive list operations.  Equality is coefficientwise equality,
checked by a finite fold.

The generic matrix layer contains dimensions, entry lookup, matrix products,
permutation or fraction-free determinants, minors, and polynomial matrices.
Cayley--Hamilton is best stated as a uniform finite identity for the generic
matrix, then specialized by coefficient substitution.  This avoids analytic
or field assumptions.  Trace identities, adjugate formulas, and finite linear
recurrences can use the same coefficientwise kernel.

Factorizations and radical expressions are certificate-relative unless a
producer has been formalized.  A factorization certificate consists of a
finite factor list, degree and leading-coefficient data, and a coefficientwise
product identity.  Its verifier is small and uniform.  Irreducibility and
Galois-theoretic impossibility results need finite permutation/group
certificates plus a fully specified arithmetization; a semantic field-theory
proof in ambient Lean is not yet an iPRA derivation.

\subsection{Roots, branches, and algebraic functions}

For approximate roots, the stable primary output is one of:

\begin{enumerate}
\item a finite multiset of isolating boxes with multiplicities;
\item a root cluster whose residual and diameter meet a requested budget;
\item a relational family of approximate roots; or
\item an individual root name under a separation or simple-root certificate.
\end{enumerate}

Ordering roots is deferred until a positive separation witness is available.
Local implicit or inverse functions consume a derivative gap, a rational box,
and a contraction or majorant certificate.  Puiseux expansions from exact
algebraic coefficient codes may use finite Newton-polygon and splitting
certificates.  The stronger demand for exact Puiseux coefficients from
arbitrary Cauchy coefficient names is not a library goal; the repository
contains a computability obstruction even for a fixed evaluator family.

\subsection{Number theory and combinatorics}

Divisibility, primality below a bound, valuations, finite divisor sums,
M\"obius inversion, totients, von Mangoldt weights, binomial coefficients,
Bell recurrences, and congruences are all exact finite arithmetic.  Their
preferred proofs use Boolean divisibility predicates and double-counting
bijections between bounded coded sets.

Asymptotic statements are expressed by explicit inequalities with a free
numerical parameter and a displayed threshold function.  A proof of
\(f(n)\sim g(n)\) becomes
\[
 \forall r\;\forall n\ge N(r)\quad
 |f(n)-g(n)|\le 2^{-r}|g(n)|,
\]
with a suitable integer formulation after clearing denominators.  The
threshold \(N\) must be primitive recursive or supplied as a certificate.
Divergence is stated by a primitive-recursive witness function: for every
height \(K\), a displayed cutoff makes the finite partial sum exceed \(K\).

Analytic number theory should separate a uniform finite verifier from the
producer of difficult constants.  Zero-free regions, exceptional-zero data,
contour bounds, and explicit estimates can be represented by rational tables
whose checker verifies all local inequalities.  This does not solve the
problem of producing effective constants in results where current proofs are
ineffective; it makes the residual obligation exact.

\subsection{Diophantine geometry}

A rational point is encoded by a normalized integer tuple with a height.
Polynomial satisfaction is decidable after clearing denominators.  Once a
height bound \(H\) is supplied, complete enumeration is a bounded search.
Conversely, from a complete finite list one obtains a height bound by taking a
finite maximum.

\begin{theorem}[Height bound--complete list equivalence]
\label{thm:height-list}
For every fixed primitive-recursive point predicate \(P(p,x)\) and height
function \(h(x)\), iPRA proves:

\begin{enumerate}
\item a bound \(H\) satisfying \(P(p,x)\to h(x)\le H\) yields a complete
      finite list of all points by bounded enumeration; and
\item a complete finite list yields the height bound equal to the maximum
      height of its members.
\end{enumerate}
\end{theorem}

\begin{proof}
For the first direction, enumerate every normalized code of height at most
\(H\), retain exactly those passing the Boolean predicate, and prove
completeness by specializing the bound hypothesis.  For the second, compute
the finite maximum.  If \(P(p,x)\), list completeness gives an index at which
\(x\) occurs, so its height is below the maximum.  Both directions are finite
folds and bounded searches.
\end{proof}

This theorem sharpens the boundary in the Faltings, Siegel, and Mordell--Weil
packets.  A verifier and enumerator can be completed independently of the
deep theorem that produces the height bound, finite cover, or rational
covering representatives.  A cardinality bound alone supplies neither.

Descent arguments should be packaged as a well-founded numerical measure,
a primitive-recursive descent step, and a proof that every nonterminal state
strictly decreases.  Primitive-recursive descent then yields the terminal
certificate.  Covering descents are relational when several next coverings
are possible; backward viability handles finite compatibility constraints.

\subsection{Geometry}

Coordinate geometry is naturally finite.  Points, lines, circles, conics, and
projective incidences are coded by rational or represented coordinates.
Determinantal incidence identities, Ceva--Menelaus-type products, power of a
point, and many triangle formulas reduce to polynomial identities plus
nondegeneracy guards.

Angle-based theorems require more care.  Prefer dot products, determinants,
complex phases, quaternions, or sine/cosine pairs to a globally chosen angle.
Normalization needs a positive norm gap.  For rotations, exact quaternion data
can remain meaningful when the vector part vanishes, whereas a unit axis at
the identity is discontinuous.  Approximate axes may be returned relationally;
exact axes require separation from the identity or a supplied nonzero vector.

Spherical and trigonometric identities should be split into exact polynomial
identities in sine/cosine variables and analytic existence/branch statements.
The first belong to the algebraic checker.  The second consume quadrant,
nonvanishing, or interval certificates.  This avoids importing classical
inverse-function conventions into the algebraic part.

\subsection{Topology}

Abstract compactness is replaced by finite combinatorics.  A quantitative
Jordan or invariance-of-dimension theorem should specify:

\begin{itemize}
\item a rational grid scale obtained from a continuity and separation modulus;
\item a finite cell or cubical complex;
\item a Boolean labeling rule;
\item a local parity or Sperner invariant; and
\item a finite path, chain, or parity certificate producing the conclusion.
\end{itemize}

The global proof is quantifier-free induction over cells or a finite trace.
Refinement from one grid to the next is a relational presentation; a coherent
geometric point is extracted only when nested cells shrink quantitatively.

\subsection{Analysis and special functions}

Calculus, Gamma and zeta identities, residues, and asymptotics should be
organized around finite sums, rational contours or meshes, and explicit
remainders.  A complex contour is a finite polygonal path or a fixed evaluator
with a modulus.  Cauchy and residue statements become finite telescoping
identities plus bounds on local approximation errors.  Winding and zero counts
are exact only under a boundary-apartness certificate.

Special functions are best given by a fixed family of finite approximants and
one or more interchangeable certificates: series tails, product tails,
integral quadrature errors, or functional-equation residuals.  Identities are
proved at finite precision through whichever representation has the easiest
error ledger.  Extensionality then shows that the alternatives describe the
same represented value.

\subsection{Probability}

The base probability library consists of finite rational laws, products,
convolutions, moments, and rational test functions.  Limit theorems are
quantitative approximation theorems.  Tail truncation, lattice spacing,
variance lower bounds, third-moment bounds, and anti-concentration are explicit
inputs or derived certificates.  Multivariate laws use finite rational vectors
and finite nets of test directions.  Martingale statements use bounded finite
filtrations and transition tables.

This formulation captures the effective finite content of central-limit
arguments while avoiding an unsupported higher-type measure object.  A more
general represented measure can be added later as an FRP whose challenges are
finite rational test functions and whose answers are certified approximate
integrals.

\subsection{A uniform package boundary}

Every domain theorem should be split into four files or namespaces:

\begin{enumerate}
\item \emph{Data}: codes, guards, evaluators, candidate bounds, and challenges;
\item \emph{Finite kernel}: Boolean folds, arithmetic identities, local
      invariants, and certificate checking;
\item \emph{Producer}: optional algorithm constructing the certificate or
      viable relation; and
\item \emph{Export}: the exact \(\Prf\,\PRAAx\) theorem and its transport by
      \code{pra\_iff\_native}.
\end{enumerate}

A packet may honestly complete the finite kernel while leaving the producer
open.  This is materially stronger and more reusable than an all-or-nothing
outline, provided the status is stated explicitly.
